% =============================================================
% File: chapters/ta/chapter-3/rancangan.tex
%
% PERUBAHAN:
%   - "Audit Log Service (ALS)" → "Audit Log Service (ALS)"
%   - Semua "ALS" → "ALS"
%   - Caption figure: "Audit Log" → "Audit Log"
%   - "rekam jejak aktivitas" → "aktivitas"
%   - "immutable" pada klaim lokal → "tamper-evident"
%     (tetap "immutable trust anchor" untuk IOTA — itu akurat)
%   - Paragraf satu baris
% =============================================================

\section{Rancangan}

Berdasarkan hasil analisis solusi pada Subbab~\ref{sec:analisis-solusi}, rancangan solusi menjelaskan bentuk sistem yang dibangun. Fondasi DecMed, yaitu IOTA sebagai penyimpanan \textit{on-chain}, IPFS sebagai penyimpanan \textit{off-chain}, dan PRE sebagai mekanisme kriptografis pendukung distribusi akses, tetap dipertahankan.

\subsection{Rancangan Arsitektur Solusi}
\label{subsec:rancangan-arsitektur}

Keseluruhan arsitektur kerangka kerja DecMed dipertahankan pada arsitektur solusi. Tiga \textit{client} utama DecMed, yaitu \textit{Patient Client}, \textit{Hospital Client}, dan \textit{Ministry Client}, tetap ada namun dilengkapi dengan instrumentasi pengiriman \textit{audit event}. PRE Server tetap berperan sebagai \textit{middleware} keamanan, dengan tambahan instrumentasi serupa.

Perubahan utama berupa penambahan satu komponen baru, yaitu \textbf{Audit Log Service (ALS)}, sebagai layanan \textit{backend} independen yang mengumpulkan, memvalidasi, dan menyimpan bukti digital dari seluruh komponen DecMed. Penambahan ALS tidak mengubah alur bisnis maupun fungsionalitas utama DecMed yang telah berjalan pada komponen eksisting. ALS bertindak sebagai pengumpul log terpusat yang memanfaatkan kembali (\textit{reuse}) infrastruktur \textit{off-chain} (IPFS Cluster) dan \textit{on-chain} (jaringan IOTA melalui Gas Station Server) yang sudah tersedia pada DecMed.

Gambaran arsitektur sistem secara keseluruhan dapat dilihat pada Gambar~\ref{fig:arsitektur-solusi}. Warna merah merepresentasikan komponen baru dan penyesuaian alur interaksi yang ditambahkan pada arsitektur DecMed.

\begin{figure}[htbp]
\centering
\includegraphics[width=1\textwidth]{chapter-3/decmed_ats_architecture.png}
\caption{Rancangan Arsitektur Solusi Audit Log DecMed}
\label{fig:arsitektur-solusi}
\end{figure}

Gambar~\ref{fig:arsitektur-solusi} menegaskan bahwa kerangka DecMed dipertahankan, sementara penambahan (ditandai merah) mencakup: (1) modul pengiriman \textit{signed event} pada setiap \textit{event source}, (2) ALS sebagai layanan pengumpul terpusat, (3) mekanisme pengarsipan berkas log ke IPFS, dan (4) publikasi metadata rotasi log ke jaringan IOTA melalui Gas Station Server yang telah ada. Dengan demikian, ALS tidak memerlukan infrastruktur penyimpanan baru, melainkan memanfaatkan ulang infrastruktur yang sudah ada pada DecMed.

Berdasarkan Gambar~\ref{fig:arsitektur-solusi}, alur interaksi ALS dalam arsitektur DecMed dirancang sebagai berikut.

\begin{enumerate}
    \item \textbf{Pengiriman \textit{event} dari komponen DecMed.} Setiap kali terjadi aktivitas bernilai audit pada \textit{event source} (yaitu \textit{Patient Client}, \textit{Hospital Client}, \textit{Ministry Client}, dan PRE Server), komponen tersebut membungkus rekaman \textit{audit event} dalam format \textit{signed payload} Ed25519 dan mengirimkannya secara asinkron ke ALS.
    \item \textbf{Pemrosesan dan pembentukan rantai integritas.} ALS menerima \textit{event}, memverifikasi tanda tangan digital pengirim, membentuk \textit{audit record} sesuai skema terstruktur, dan menerapkan mekanisme \textit{hash-chaining} antar \textit{record} yang masuk.
    \item \textbf{Penyimpanan berkas \textit{off-chain} di IPFS.} Secara periodik, ALS merotasi berkas log aktif, mengunggah berkas arsip ke IPFS Cluster, dan menerapkan \textit{pinning} agar berkas tidak terhapus oleh proses \textit{Garbage Collection}.
    \item \textbf{Pencatatan metadata \textit{on-chain} di IOTA.} ALS menerbitkan transaksi ke jaringan IOTA melalui Gas Station Server untuk mencatatkan metadata rotasi log ke dalam \textit{smart contract} Move sebagai \textit{immutable trust anchor}.
\end{enumerate}

Dengan rancangan arsitektur ini, ALS berhasil memisahkan beban pemrosesan dan penyimpanan log dari transaksi data klinis utama, sekaligus menjamin bahwa seluruh aktivitas pada DecMed terekam secara utuh, bersifat \textit{tamper-evident} pada level lokal, dan \textit{immutable} setelah tersimpan pada IOTA dan IPFS.

\subsection{Rancangan Skema Audit Record}
\label{subsec:rancangan-audit-record}

\subsubsection{Atribut Umum Audit Record}

Berdasarkan pemetaan ancaman terhadap audit event pada Tabel~\ref{tab:threat-event-mapping}, setiap audit event memiliki kebutuhan bukti yang berbeda sesuai dengan aktivitas dan ancaman yang direpresentasikan. Meskipun demikian, terdapat beberapa informasi yang memiliki fungsi umum dan dapat digunakan pada seluruh jenis audit event. Informasi tersebut diperlukan agar setiap audit record dapat menjelaskan pihak yang terlibat, waktu terjadinya aktivitas, objek yang berkaitan dengan aktivitas, serta hasil dari aktivitas tersebut.

Empat kebutuhan bukti umum yang diidentifikasi adalah \textit{actor}, \textit{timestamp}, \textit{object}, dan \textit{result}. \textit{Actor} digunakan untuk mengidentifikasi pihak atau komponen yang melakukan atau memicu aktivitas yang dicatat. Identitas tersebut tidak selalu berupa identitas pengguna, karena beberapa aktivitas dapat dilakukan oleh komponen sistem secara otomatis. Oleh karena itu, atribut ini dapat merepresentasikan pengguna, administrator, proses, maupun komponen sistem yang terkait dengan suatu aktivitas.

\textit{Timestamp} digunakan untuk mencatat waktu terjadinya aktivitas. Informasi waktu diperlukan untuk mengetahui urutan kejadian, melakukan korelasi antar-event, serta menganalisis pola aktivitas yang berkaitan dengan suatu ancaman. Pencatatan waktu juga memungkinkan audit record dikaitkan dengan kejadian lain yang terjadi pada periode yang sama.

\textit{Object} digunakan untuk mengidentifikasi sumber daya atau entitas yang menjadi sasaran aktivitas. Bentuk objek bergantung pada jenis event, seperti rekam medis, \textit{capability}, transaksi IOTA, objek IPFS, objek Redis, maupun metadata \textit{on-chain}. Dengan menggunakan atribut umum \textit{object}, perbedaan jenis sumber daya tersebut dapat direpresentasikan tanpa harus menjadikan setiap jenis sumber daya sebagai atribut pada seluruh audit record.

Sementara itu, \textit{result} digunakan untuk mencatat hasil dari aktivitas yang dilakukan. Hasil dapat berupa keberhasilan, kegagalan, penolakan, maupun status lain yang relevan dengan operasi. Informasi ini diperlukan untuk membedakan aktivitas yang berhasil dilakukan dari percobaan aktivitas yang gagal atau ditolak, sehingga dapat digunakan dalam investigasi maupun analisis ancaman.

Keempat kebutuhan bukti tersebut kemudian dipetakan menjadi atribut umum pada struktur audit record, sebagaimana ditunjukkan pada Tabel~\ref{tab:general-audit-evidence}. Atribut lain yang hanya relevan terhadap jenis event tertentu tetap direpresentasikan sebagai atribut khusus pada masing-masing event.

\begin{xltabular}
    {\textwidth}
    {|>{\raggedright\arraybackslash}p{3.2cm}
     |>{\centering\arraybackslash}p{2.8cm}
     |X|}
    
\caption{Pemetaan Kebutuhan Bukti Umum ke Atribut Audit Record}
\label{tab:general-audit-evidence}\\

\hline
\textbf{Kebutuhan Bukti} &
\textbf{Atribut Umum} &
\textbf{Keterangan} \\
\hline
\endfirsthead

\hline
\textbf{Kebutuhan Bukti} &
\textbf{Atribut Umum} &
\textbf{Keterangan} \\
\hline
\endhead

\hline
\endfoot

\hline
\endlastfoot

Identitas pihak atau komponen yang melakukan atau memicu aktivitas &
\texttt{actor} &
Mengidentifikasi pengguna, administrator, proses, atau komponen sistem yang melakukan atau memicu aktivitas audit. \\

\hline

Waktu terjadinya aktivitas &
\texttt{timestamp} &
Mencatat waktu terjadinya aktivitas untuk menentukan urutan kejadian dan melakukan korelasi antar-event. \\

\hline

Objek atau sumber daya yang menjadi sasaran aktivitas &
\texttt{object} &
Mengidentifikasi entitas yang berkaitan dengan aktivitas, seperti rekam medis, \textit{capability}, transaksi, objek IPFS, objek Redis, atau metadata \textit{on-chain}. \\

\hline

Hasil aktivitas atau pemrosesan &
\texttt{result} &
Mencatat hasil aktivitas, seperti berhasil, gagal, ditolak, atau status lain yang relevan. \\

\hline

\end{xltabular}

\subsubsection{Atribut Khusus Audit Record}

Selain atribut umum tersebut, setiap \textit{audit event} juga memiliki atribut tambahan yang disesuaikan dengan kebutuhan bukti digital terhadap ancaman yang dipetakan pada Tabel\ref{tab:threat-event-mapping}. Atribut tambahan tersebut disajikan pada Tabel~\ref{tab:detail-audit-record}.

\subsection{Rancangan Spesifikasi \textit{Audit Event}}
\label{sec:spesifikasi-audit-event}



Subbab ini merinci spesifikasi setiap \textit{audit event} yang telah
diidentifikasi pada Tabel~\ref{tab:threat-event-mapping}. Untuk setiap
\textit{event}, dijelaskan konteks aktivitas yang dicatat, kemungkinan nilai
pada atribut umum, serta atribut tambahan beserta kemungkinan nilainya.

Atribut umum yang terdapat pada seluruh jenis \textit{audit event} adalah
\texttt{source\_component}, \texttt{actor\_id}, \texttt{target\_object},
\texttt{outcome}, dan \texttt{action\_type}. Kemungkinan nilai dari atribut
\texttt{outcome} dan \texttt{action\_type} berlaku seragam di seluruh
\textit{event} dan dijelaskan pada Tabel~\ref{tab:nilai-outcome} dan
Tabel~\ref{tab:nilai-action-type}.

\begin{table}[htbp]
\centering
\caption{Kemungkinan Nilai Atribut \texttt{outcome}}
\label{tab:nilai-outcome}
\renewcommand{\arraystretch}{1.3}
\begin{tabular}{|p{2.5cm}|p{9.5cm}|}
\hline
\textbf{Nilai} & \textbf{Kondisi} \\
\hline
\texttt{Success} &
Aktivitas berhasil diselesaikan sesuai dengan yang dimaksudkan. \\
\hline
\texttt{Failure} &
Aktivitas gagal diselesaikan akibat kesalahan teknis, kondisi yang tidak
terpenuhi, atau kegagalan pada komponen yang dihubungi (mis. IOTA, IPFS,
Redis tidak merespons). \\
\hline
\texttt{Denied} &
Aktivitas ditolak secara eksplisit karena tidak memenuhi persyaratan
otorisasi, misalnya \textit{capability} tidak valid, JWT kadaluarsa, atau
peran tidak sesuai. \\
\hline
\end{tabular}
\end{table}

\begin{table}[htbp]
\centering
\caption{Kemungkinan Nilai Atribut \texttt{action\_type}}
\label{tab:nilai-action-type}
\renewcommand{\arraystretch}{1.3}
\begin{tabular}{|p{2.5cm}|p{9.5cm}|}
\hline
\textbf{Nilai} & \textbf{Kondisi} \\
\hline
\texttt{CREATE} &
Membuat atau menginisialisasi objek, entri, atau sumber daya baru.
Contoh: mendaftarkan akun, menerbitkan \textit{activation key}, membuat
entri rekam medis, mengunggah konten ke IPFS. \\
\hline
\texttt{READ} &
Membaca atau mengambil data dari objek yang sudah ada tanpa melakukan
perubahan. Contoh: membaca rekam medis, mengambil data administratif,
mengunduh konten dari IPFS. \\
\hline
\texttt{UPDATE} &
Memperbarui atau memodifikasi objek yang sudah ada. Contoh: memperbarui
entri rekam medis, memperbarui \textit{activation key}, memperbarui
metadata administratif. \\
\hline
\texttt{DELETE} &
Menghapus atau mencabut objek, entri, atau izin yang sudah ada. Contoh:
mencabut delegasi akses (\texttt{revoke\_access}), \textit{sign out},
menghapus \textit{nonce} dari Redis setelah digunakan. \\
\hline
\texttt{VALIDATE} &
Memverifikasi keabsahan, keaslian, atau kelengkapan suatu objek atau
klaim tanpa mengubah \textit{state}. Contoh: memvalidasi QR \textit{payload},
memvalidasi JWT, memvalidasi \textit{capability on-chain}. \\
\hline
\texttt{EXECUTE} &
Menjalankan proses yang telah dipersiapkan sebelumnya. Berbeda dari
\texttt{CREATE} karena objeknya sudah ada dan tinggal dieksekusi. Contoh:
mengeksekusi transaksi yang telah disponsori \textit{gas}. \\
\hline
\end{tabular}
\end{table}

% -----------------------------------------------------------------------
\subsubsection{EV1 --- Authentication}
\label{subsec:ev1-authentication}

EV1 mencatat setiap percobaan autentikasi pengguna pada \textit{Patient Client}
maupun \textit{Hospital Client}, baik yang berhasil maupun yang gagal. Aktivitas
yang tercakup meliputi proses \textit{sign-in} menggunakan \textit{activation key},
\textit{patient ID}, atau \textit{seed words}, serta proses \textit{sign-out}.
Pencatatan \textit{event} ini diperlukan untuk mendeteksi pola percobaan autentikasi
yang mencurigakan, seperti percobaan berulang yang mengindikasikan serangan
\textit{brute force}.

\begin{table}[htbp]
\centering
\caption{Atribut Umum EV1 Authentication}
\label{tab:ev1-umum}
\renewcommand{\arraystretch}{1.3}
\begin{tabular}{|p{3.5cm}|p{8.5cm}|}
\hline
\textbf{Atribut} & \textbf{Kemungkinan Nilai} \\
\hline
\texttt{source\_component} &
\texttt{HospitalTauriBackend} | \texttt{PatientTauriBackend} \\
\hline
\texttt{actor\_id} &
IOTA \textit{address} pengguna yang melakukan autentikasi. \\
\hline
\texttt{target\_object} &
\texttt{activation\_key} | \texttt{patient\_id} | \texttt{seed\_words}
--- menyatakan mekanisme autentikasi yang digunakan. \\
\hline
\texttt{action\_type} &
\texttt{VALIDATE} | \texttt{CREATE} | \texttt{DELETE} --- \texttt{VALIDATE}
untuk verifikasi kredensial, \texttt{CREATE} untuk inisialisasi sesi,
\texttt{DELETE} untuk \textit{sign-out}. \\
\hline
\texttt{outcome} &
\texttt{Success} | \texttt{Failure} | \texttt{Denied} \\
\hline
\end{tabular}
\end{table}

\begin{table}[htbp]
\centering
\caption{Atribut Tambahan EV1 Authentication}
\label{tab:ev1-detail}
\renewcommand{\arraystretch}{1.3}
\begin{tabular}{|p{4cm}|p{7.5cm}|}
\hline
\textbf{Atribut} & \textbf{Kemungkinan Nilai dan Keterangan} \\
\hline
\texttt{auth\_method} &
\texttt{activation\_key} | \texttt{patient\_id} | \texttt{seed\_words}
--- metode spesifik yang digunakan dalam percobaan autentikasi. \\
\hline
\texttt{role} &
\texttt{Pasien} | \texttt{PersonnelMedis} | \texttt{PersonnelAdministratif}
| \texttt{AdminFasyankes} --- peran pengguna yang sedang melakukan autentikasi. \\
\hline
\texttt{attempt\_count} &
Bilangan bulat positif --- jumlah percobaan autentikasi yang telah dilakukan
pada sesi atau dalam periode tertentu. \\
\hline
\texttt{failure\_reason} &
\texttt{invalid\_credentials} | \texttt{account\_not\_found} |
\texttt{session\_expired} | \texttt{null} apabila \texttt{outcome} bukan
\texttt{Failure}. \\
\hline
\texttt{device\_info} &
Informasi perangkat yang digunakan, misalnya \textit{fingerprint} perangkat
atau platform sistem operasi. \\
\hline
\end{tabular}
\end{table}

% -----------------------------------------------------------------------
\subsubsection{EV2 --- QR Validation}
\label{subsec:ev2-qr-validation}

EV2 mencatat proses validasi \textit{payload} QR Access Delegation yang dipindai
oleh pasien dari QR \textit{code} yang ditampilkan oleh \textit{hospital personnel}.
\textit{Payload} QR berisi IOTA \textit{address} dan PRE \textit{public key}
\textit{personnel} dalam format \texttt{<iota\_address>@<hex\_pre\_pub\_key>}.
Pencatatan \textit{event} ini diperlukan untuk mendeteksi pemalsuan
\textit{payload} QR maupun kebocoran informasi delegasi akses.

\begin{table}[htbp]
\centering
\caption{Atribut Umum EV2 QR Validation}
\label{tab:ev2-umum}
\renewcommand{\arraystretch}{1.3}
\begin{tabular}{|p{3.5cm}|p{8.5cm}|}
\hline
\textbf{Atribut} & \textbf{Kemungkinan Nilai} \\
\hline
\texttt{source\_component} &
\texttt{PatientTauriBackend} \\
\hline
\texttt{actor\_id} &
IOTA \textit{address} pasien yang memindai QR. \\
\hline
\texttt{target\_object} &
IOTA \textit{address} \& PRE \textit{public key hospital personnel} yang
terdapat dalam QR. \\
\hline
\texttt{action\_type} &
\texttt{READ} | \texttt{VALIDATE} \\
\hline
\texttt{outcome} &
\texttt{Success} | \texttt{Failure} \\
\hline
\end{tabular}
\end{table}

\begin{table}[htbp]
\centering
\caption{Atribut Tambahan EV2 QR Validation}
\label{tab:ev2-detail}
\renewcommand{\arraystretch}{1.3}
\begin{tabular}{|p{4cm}|p{7.5cm}|}
\hline
\textbf{Atribut} & \textbf{Kemungkinan Nilai dan Keterangan} \\
\hline
\texttt{qr\_content\_hash} &
Nilai \textit{hash} SHA-256 dari \textit{payload} QR yang dipindai ---
digunakan untuk membuktikan konten yang diterima tanpa menyimpan nilai
aslinya. \\
\hline
\texttt{hospital\_personnel\_iota\_address} &
IOTA \textit{address} \textit{hospital personnel} yang terkandung dalam
QR. \\
\hline
\texttt{hospital\_personnel\_pre\_public\_key\_fingerprint} &
\textit{Fingerprint} (hash) dari PRE \textit{public key personnel} yang
terkandung dalam QR. \\
\hline
\texttt{hospital\_name} &
Nama fasyankes tempat \textit{personnel} bernaung, apabila tersedia dalam
\textit{payload}. \\
\hline
\texttt{hospital\_personnel\_name} &
Nama \textit{personnel} yang ditampilkan dalam QR, apabila tersedia. \\
\hline
\texttt{validation\_step} &
\texttt{parse} | \texttt{format\_check} | \texttt{address\_lookup} ---
tahapan validasi terakhir yang berhasil atau gagal. \\
\hline
\end{tabular}
\end{table}

% -----------------------------------------------------------------------
\subsubsection{EV3 --- Capability Creation}
\label{subsec:ev3-capability-creation}

EV3 mencatat proses pembuatan dan pencabutan \textit{capability} akses yang
diberikan pasien kepada \textit{hospital personnel}. \textit{Capability} ini
tersimpan sebagai entri \textit{on-chain} (HospitalPersonnelAccess) dan \textit{kfrag}
di Redis. Pencatatan \textit{event} ini diperlukan untuk membuktikan bahwa
pasien benar-benar telah memberikan atau mencabut delegasi akses kepada
\textit{personnel} tertentu, sehingga menangkal ancaman penyangkalan (T6).

\begin{table}[htbp]
\centering
\caption{Atribut Umum EV3 Capability Creation}
\label{tab:ev3-umum}
\renewcommand{\arraystretch}{1.3}
\begin{tabular}{|p{3.5cm}|p{8.5cm}|}
\hline
\textbf{Atribut} & \textbf{Kemungkinan Nilai} \\
\hline
\texttt{source\_component} &
\texttt{PatientTauriBackend} | \texttt{ProxyReencryption} \\
\hline
\texttt{actor\_id} &
IOTA \textit{address} pasien yang mendelegasikan akses. \\
\hline
\texttt{target\_object} &
IOTA \textit{address} \textit{hospital personnel} penerima akses. \\
\hline
\texttt{action\_type} &
\texttt{CREATE} | \texttt{DELETE} --- \texttt{CREATE} untuk pemberian
akses, \texttt{DELETE} untuk pencabutan akses. \\
\hline
\texttt{outcome} &
\texttt{Success} | \texttt{Failure} \\
\hline
\end{tabular}
\end{table}

\begin{table}[htbp]
\centering
\caption{Atribut Tambahan EV3 Capability Creation}
\label{tab:ev3-detail}
\renewcommand{\arraystretch}{1.3}
\begin{tabular}{|p{4cm}|p{7.5cm}|}
\hline
\textbf{Atribut} & \textbf{Kemungkinan Nilai dan Keterangan} \\
\hline
\texttt{access\_type} &
\texttt{read} | \texttt{update} --- cakupan akses yang diberikan. \\
\hline
\texttt{exp\_duration\_ms} &
Durasi kedaluwarsa \textit{capability} dalam milidetik. \\
\hline
\texttt{k\_frag\_fingerprint} &
\textit{Fingerprint} (hash) dari \textit{kfrag} yang disimpan di Redis,
tanpa mengekspos nilai \textit{kfrag} asli. \\
\hline
\texttt{transaction\_digest} &
\textit{Digest} transaksi IOTA yang merekam pembuatan atau pencabutan
\textit{capability on-chain}. \\
\hline
\texttt{hospital\_name} &
Nama fasyankes tempat \textit{personnel} penerima akses bernaung. \\
\hline
\texttt{nonce\_used} &
\textit{Nonce} yang digunakan dalam proses pemberian akses untuk mencegah
\textit{replay attack}. \\
\hline
\end{tabular}
\end{table}

% -----------------------------------------------------------------------
\subsubsection{EV4 --- Medical Record Access}
\label{subsec:ev4-medical-record-access}

EV4 mencatat setiap aktivitas pembacaan maupun penulisan data rekam medis
yang dieksekusi oleh PRE Server atas permintaan \textit{hospital personnel}.
Aktivitas ini mencakup pengambilan data administratif, pengambilan data klinis
beserta proses \textit{re-encryption}, pembuatan entri rekam medis baru, dan
pembaruan entri rekam medis. Pencatatan \textit{event} ini diperlukan untuk
membuktikan bahwa akses terhadap data sensitif pasien dilakukan dengan otorisasi
yang sah (T7, T13).

\begin{table}[htbp]
\centering
\caption{Atribut Umum EV4 Medical Record Access}
\label{tab:ev4-umum}
\renewcommand{\arraystretch}{1.3}
\begin{tabular}{|p{3.5cm}|p{8.5cm}|}
\hline
\textbf{Atribut} & \textbf{Kemungkinan Nilai} \\
\hline
\texttt{source\_component} &
\texttt{ProxyReencryption} \\
\hline
\texttt{actor\_id} &
IOTA \textit{address} \textit{hospital personnel} yang mengakses. \\
\hline
\texttt{target\_object} &
IPFS CID rekam medis (untuk data klinis) | IOTA \textit{address} pasien
(untuk data administratif). \\
\hline
\texttt{action\_type} &
\texttt{CREATE} | \texttt{READ} | \texttt{UPDATE} \\
\hline
\texttt{outcome} &
\texttt{Success} | \texttt{Failure} | \texttt{Denied} \\
\hline
\end{tabular}
\end{table}

\begin{table}[htbp]
\centering
\caption{Atribut Tambahan EV4 Medical Record Access}
\label{tab:ev4-detail}
\renewcommand{\arraystretch}{1.3}
\begin{tabular}{|p{4cm}|p{7.5cm}|}
\hline
\textbf{Atribut} & \textbf{Kemungkinan Nilai dan Keterangan} \\
\hline
\texttt{patient\_iota\_address} &
IOTA \textit{address} pasien pemilik data yang diakses. \\
\hline
\texttt{record\_index} &
Indeks rekam medis yang diakses atau dibuat, apabila relevan. \\
\hline
\texttt{role\_used} &
\texttt{MedicalPersonnel} | \texttt{AdministrativePersonnel} --- peran
\textit{personnel} yang tercantum pada JWT saat akses dilakukan. \\
\hline
\texttt{purpose\_used} &
\texttt{read} | \texttt{update} --- tujuan akses yang tercantum pada JWT. \\
\hline
\texttt{jwt\_sub} &
\textit{Subject} JWT yang digunakan --- biasanya IOTA \textit{address
personnel}. \\
\hline
\texttt{capability\_valid} &
\texttt{true} | \texttt{false} --- hasil validasi \textit{capability
on-chain} sebelum akses diproses. \\
\hline
\texttt{ipfs\_cid} &
CID IPFS konten klinis yang diakses atau diunggah, \texttt{null} apabila
operasi hanya menyangkut data administratif. \\
\hline
\texttt{reencryption\_performed} &
\texttt{true} | \texttt{false} --- menyatakan apakah proses
\textit{re-encryption} PRE dilakukan pada operasi ini. \\
\hline
\end{tabular}
\end{table}

% -----------------------------------------------------------------------
\subsubsection{EV5 --- Personnel Activation Key}
\label{subsec:ev5-personnel-activation-key}

EV5 mencatat aktivitas penerbitan dan pembaruan \textit{activation key} yang
dilakukan oleh Admin Fasyankes untuk \textit{personnel} di bawahnya.
\textit{Activation key} ini memungkinkan \textit{personnel} mendaftarkan diri
ke dalam sistem. Pencatatan \textit{event} ini diperlukan untuk membuktikan
bahwa admin fasyankes benar-benar menerbitkan \textit{activation key} kepada
\textit{personnel} tertentu (T8).

\begin{table}[htbp]
\centering
\caption{Atribut Umum EV5 Personnel Activation Key}
\label{tab:ev5-umum}
\renewcommand{\arraystretch}{1.3}
\begin{tabular}{|p{3.5cm}|p{8.5cm}|}
\hline
\textbf{Atribut} & \textbf{Kemungkinan Nilai} \\
\hline
\texttt{source\_component} &
\texttt{HospitalTauriBackend} \\
\hline
\texttt{actor\_id} &
IOTA \textit{address} Admin Fasyankes yang menerbitkan \textit{key}. \\
\hline
\texttt{target\_object} &
\textit{Hash} ID personel penerima \textit{activation key}. \\
\hline
\texttt{action\_type} &
\texttt{CREATE} | \texttt{UPDATE} \\
\hline
\texttt{outcome} &
\texttt{Success} | \texttt{Failure} \\
\hline
\end{tabular}
\end{table}

\begin{table}[htbp]
\centering
\caption{Atribut Tambahan EV5 Personnel Activation Key}
\label{tab:ev5-detail}
\renewcommand{\arraystretch}{1.3}
\begin{tabular}{|p{4cm}|p{7.5cm}|}
\hline
\textbf{Atribut} & \textbf{Kemungkinan Nilai dan Keterangan} \\
\hline
\texttt{admin\_iota\_address} &
IOTA \textit{address} Admin Fasyankes yang menerbitkan \textit{key}. \\
\hline
\texttt{personnel\_id\_hash} &
\textit{Hash} dari ID personel penerima \textit{activation key}. \\
\hline
\texttt{role\_assigned} &
\texttt{MedicalPersonnel} | \texttt{AdministrativePersonnel} | \texttt{Admin}
--- peran yang ditetapkan pada \textit{activation key}. \\
\hline
\texttt{hospital\_id\_hash} &
\textit{Hash} dari ID fasyankes tempat \textit{personnel} bernaung. \\
\hline
\texttt{transaction\_digest} &
\textit{Digest} transaksi IOTA yang merekam penerbitan atau pembaruan
\textit{activation key on-chain}. \\
\hline
\texttt{activation\_key\_hash} &
\textit{Hash} dari nilai \textit{activation key} yang diterbitkan. \\
\hline
\end{tabular}
\end{table}

% -----------------------------------------------------------------------
\subsubsection{EV6 --- Facility Registration}
\label{subsec:ev6-facility-registration}

EV6 mencatat aktivitas registrasi fasyankes baru beserta Admin Fasyankes-nya
yang dilakukan oleh Ministry Administrator. Pencatatan \textit{event} ini
diperlukan untuk membuktikan bahwa Kementerian benar-benar telah meregistrasi
fasyankes tertentu beserta administratornya (T9).

\begin{table}[htbp]
\centering
\caption{Atribut Umum EV6 Facility Registration}
\label{tab:ev6-umum}
\renewcommand{\arraystretch}{1.3}
\begin{tabular}{|p{3.5cm}|p{8.5cm}|}
\hline
\textbf{Atribut} & \textbf{Kemungkinan Nilai} \\
\hline
\texttt{source\_component} &
\texttt{MinistryTauriBackend} \\
\hline
\texttt{actor\_id} &
IOTA \textit{address} Ministry Administrator. \\
\hline
\texttt{target\_object} &
\textit{Hash} ID fasyankes yang diregistrasi. \\
\hline
\texttt{action\_type} &
\texttt{CREATE} | \texttt{UPDATE} \\
\hline
\texttt{outcome} &
\texttt{Success} | \texttt{Failure} \\
\hline
\end{tabular}
\end{table}

\begin{table}[htbp]
\centering
\caption{Atribut Tambahan EV6 Facility Registration}
\label{tab:ev6-detail}
\renewcommand{\arraystretch}{1.3}
\begin{tabular}{|p{4cm}|p{7.5cm}|}
\hline
\textbf{Atribut} & \textbf{Kemungkinan Nilai dan Keterangan} \\
\hline
\texttt{facility\_id} &
Identitas fasyankes yang diregistrasi. \\
\hline
\texttt{facility\_name} &
Nama fasyankes yang diregistrasi. \\
\hline
\texttt{administrator\_id} &
Identitas Admin Fasyankes yang ditetapkan saat registrasi. \\
\hline
\texttt{transaction\_digest} &
\textit{Digest} transaksi IOTA yang merekam registrasi fasyankes. \\
\hline
\texttt{activation\_key\_hash} &
\textit{Hash} dari \textit{facility activation key} yang diterbitkan
saat registrasi. \\
\hline
\end{tabular}
\end{table}

% -----------------------------------------------------------------------
\subsubsection{EV7 --- PRE Service Operation}
\label{subsec:ev7-pre-service-operation}

EV7 mencatat setiap permintaan layanan yang diterima oleh PRE Server dari
komponen \textit{client} maupun dari PRE Server itu sendiri. Aktivitas yang
tercakup meliputi seluruh \textit{endpoint} REST PRE Server: \texttt{/nonce},
\texttt{/keys}, \texttt{/medical-record}, dan \texttt{/administrative}.
Pencatatan \textit{event} ini diperlukan untuk mendeteksi penyamaran sebagai
PRE Server (\textit{man-in-the-middle}), penyadapan komunikasi, dan pembanjiran
\textit{endpoint} (T10, T11, T12).

\begin{table}[htbp]
\centering
\caption{Atribut Umum EV7 PRE Service Operation}
\label{tab:ev7-umum}
\renewcommand{\arraystretch}{1.3}
\begin{tabular}{|p{3.5cm}|p{8.5cm}|}
\hline
\textbf{Atribut} & \textbf{Kemungkinan Nilai} \\
\hline
\texttt{source\_component} &
\texttt{ProxyReencryption} \\
\hline
\texttt{actor\_id} &
IOTA \textit{address} pemohon | \texttt{proxy\_iota\_address} apabila
PRE Server bertindak sebagai aktor. \\
\hline
\texttt{target\_object} &
\textit{Endpoint} PRE Server yang dipanggil, misalnya
\texttt{/api/v1/nonce}, \texttt{/api/v1/keys},
\texttt{/api/v1/medical-record}. \\
\hline
\texttt{action\_type} &
\texttt{CREATE} | \texttt{READ} | \texttt{VALIDATE} \\
\hline
\texttt{outcome} &
\texttt{Success} | \texttt{Failure} | \texttt{Denied} \\
\hline
\end{tabular}
\end{table}

\begin{table}[htbp]
\centering
\caption{Atribut Tambahan EV7 PRE Service Operation}
\label{tab:ev7-detail}
\renewcommand{\arraystretch}{1.3}
\begin{tabular}{|p{4cm}|p{7.5cm}|}
\hline
\textbf{Atribut} & \textbf{Kemungkinan Nilai dan Keterangan} \\
\hline
\texttt{endpoint\_called} &
\textit{Path} lengkap \textit{endpoint} yang dipanggil, misalnya
\texttt{/api/v1/medical-record}. \\
\hline
\texttt{request\_id} &
Pengenal unik permintaan yang dapat digunakan untuk menghubungkan
\textit{event} ini dengan \textit{event} lain pada siklus yang sama. \\
\hline
\texttt{caller\_component} &
\texttt{PatientTauriBackend} | \texttt{HospitalTauriBackend} |
\texttt{ProxyReencryption} --- komponen yang mengajukan permintaan. \\
\hline
\texttt{channel\_encryption} &
\texttt{TLS} | \texttt{none} --- mekanisme pengamanan kanal komunikasi
yang digunakan. \\
\hline
\texttt{jwt\_purpose} &
\texttt{read} | \texttt{update} | \texttt{null} --- tujuan yang tercantum
pada JWT yang disertakan dalam permintaan, apabila ada. \\
\hline
\texttt{http\_status\_code} &
Kode status HTTP yang dikembalikan PRE Server sebagai respons, misalnya
\texttt{200}, \texttt{401}, \texttt{403}, \texttt{500}. \\
\hline
\texttt{latency\_ms} &
Durasi pemrosesan permintaan oleh PRE Server dalam milidetik. \\
\hline
\end{tabular}
\end{table}

% -----------------------------------------------------------------------
\subsubsection{EV8 --- IOTA Transaction Submission}
\label{subsec:ev8-iota-transaction}

EV8 mencatat setiap pengiriman transaksi Move ke jaringan IOTA, baik yang
dilakukan oleh komponen \textit{client} (Patient, Hospital, Ministry) maupun
oleh PRE Server secara langsung menggunakan \texttt{proxy\_iota\_address}-nya
sendiri. Pencatatan \textit{event} ini diperlukan untuk mendeteksi modifikasi
transaksi selama pengiriman (T14).

\begin{table}[htbp]
\centering
\caption{Atribut Umum EV8 IOTA Transaction Submission}
\label{tab:ev8-umum}
\renewcommand{\arraystretch}{1.3}
\begin{tabular}{|p{3.5cm}|p{8.5cm}|}
\hline
\textbf{Atribut} & \textbf{Kemungkinan Nilai} \\
\hline
\texttt{source\_component} &
\texttt{HospitalTauriBackend} | \texttt{PatientTauriBackend} |
\texttt{MinistryTauriBackend} | \texttt{ProxyReencryption} \\
\hline
\texttt{actor\_id} &
IOTA \textit{address} penandatangan transaksi | \texttt{proxy\_iota\_address}
apabila PRE Server yang mengirim. \\
\hline
\texttt{target\_object} &
Nama fungsi Move yang dipanggil, misalnya \texttt{create\_access},
\texttt{signup}, \texttt{create\_medical\_record}. \\
\hline
\texttt{action\_type} &
\texttt{CREATE} \\
\hline
\texttt{outcome} &
\texttt{Success} | \texttt{Failure} \\
\hline
\end{tabular}
\end{table}

\begin{table}[htbp]
\centering
\caption{Atribut Tambahan EV8 IOTA Transaction Submission}
\label{tab:ev8-detail}
\renewcommand{\arraystretch}{1.3}
\begin{tabular}{|p{4cm}|p{7.5cm}|}
\hline
\textbf{Atribut} & \textbf{Kemungkinan Nilai dan Keterangan} \\
\hline
\texttt{transaction\_digest} &
\textit{Digest} transaksi IOTA yang dihasilkan setelah transaksi berhasil
dikirim dan dikonfirmasi. \\
\hline
\texttt{payload\_hash} &
Nilai \textit{hash} dari \textit{payload} transaksi sebelum dikirim ---
digunakan untuk mendeteksi modifikasi selama pengiriman. \\
\hline
\texttt{move\_function} &
Nama fungsi Move yang dipanggil, misalnya \texttt{create\_access},
\texttt{revoke\_access}, \texttt{create\_medical\_record}. \\
\hline
\texttt{move\_module} &
Nama modul Move tempat fungsi berada, misalnya \texttt{patient},
\texttt{proxy}, \texttt{admin}. \\
\hline
\texttt{network\_confirmation\_status} &
\texttt{confirmed} | \texttt{pending} | \texttt{failed} --- status
konfirmasi transaksi pada jaringan IOTA. \\
\hline
\texttt{gas\_used} &
Jumlah \textit{gas} yang dikonsumsi oleh transaksi. \\
\hline
\texttt{sponsor\_address} &
IOTA \textit{address} Gas Station yang mensponsori transaksi, apabila
transaksi menggunakan mekanisme sponsorship. \\
\hline
\end{tabular}
\end{table}

% -----------------------------------------------------------------------
\subsubsection{EV9 --- Gas Sponsorship Request}
\label{subsec:ev9-gas-sponsorship}

EV9 mencatat setiap permintaan \textit{sponsorship} transaksi kepada Gas
Station, baik dari komponen \textit{client} maupun dari PRE Server. Permintaan
ini mencakup dua tahap: reservasi \textit{gas object} (\texttt{/reserve\_gas})
dan eksekusi transaksi yang telah disponsori (\texttt{/execute\_tx}).
Pencatatan \textit{event} ini diperlukan untuk mendeteksi penyamaran sebagai
pemohon sah, penyangkalan atas permintaan \textit{sponsorship}, dan
pembanjiran \textit{gas object} (T15, T16, T17).

\begin{table}[htbp]
\centering
\caption{Atribut Umum EV9 Gas Sponsorship Request}
\label{tab:ev9-umum}
\renewcommand{\arraystretch}{1.3}
\begin{tabular}{|p{3.5cm}|p{8.5cm}|}
\hline
\textbf{Atribut} & \textbf{Kemungkinan Nilai} \\
\hline
\texttt{source\_component} &
\texttt{HospitalTauriBackend} | \texttt{PatientTauriBackend} |
\texttt{MinistryTauriBackend} | \texttt{ProxyReencryption} \\
\hline
\texttt{actor\_id} &
IOTA \textit{address} pemohon \textit{sponsorship} |
\texttt{proxy\_iota\_address} apabila PRE Server yang memohon. \\
\hline
\texttt{target\_object} &
\textit{Endpoint} Gas Station yang dituju: \texttt{/reserve\_gas} |
\texttt{/execute\_tx}. \\
\hline
\texttt{action\_type} &
\texttt{CREATE} | \texttt{EXECUTE} \\
\hline
\texttt{outcome} &
\texttt{Success} | \texttt{Failure} \\
\hline
\end{tabular}
\end{table}

\begin{table}[htbp]
\centering
\caption{Atribut Tambahan EV9 Gas Sponsorship Request}
\label{tab:ev9-detail}
\renewcommand{\arraystretch}{1.3}
\begin{tabular}{|p{4cm}|p{7.5cm}|}
\hline
\textbf{Atribut} & \textbf{Kemungkinan Nilai dan Keterangan} \\
\hline
\texttt{gas\_budget\_requested} &
Jumlah \textit{gas budget} yang diminta dalam permintaan \textit{sponsorship}. \\
\hline
\texttt{reserve\_duration\_secs} &
Durasi reservasi \textit{gas object} dalam detik sebelum kadaluarsa. \\
\hline
\texttt{reservation\_id} &
Pengenal unik reservasi \textit{gas} yang diterbitkan oleh Gas Station. \\
\hline
\texttt{sponsor\_address} &
IOTA \textit{address} Gas Station yang bertindak sebagai \textit{sponsor}. \\
\hline
\texttt{transaction\_digest} &
\textit{Digest} transaksi yang disponsori, tersedia setelah tahap eksekusi. \\
\hline
\texttt{gas\_coin\_object\_ids} &
Daftar ID objek \textit{gas coin} yang direservasi atau digunakan dalam
transaksi. \\
\hline
\end{tabular}
\end{table}

% -----------------------------------------------------------------------
\subsubsection{EV10 --- Redis Operation}
\label{subsec:ev10-redis-operation}

EV10 mencatat setiap operasi penyimpanan, pembacaan, atau penghapusan objek
sensitif pada Redis yang dilakukan oleh PRE Server. Objek yang dikelola
mencakup \textit{nonce} autentikasi (kunci Redis: \texttt{nonce:<iota\_address>})
dan material PRE berupa \textit{kfrag} serta \textit{access keys} (kunci Redis:
\texttt{keys:<hp\_addr>@<patient\_addr>}). Pencatatan \textit{event} ini
diperlukan untuk mendeteksi akses tidak sah atau modifikasi objek Redis (T18,
T19, T20).

\begin{table}[htbp]
\centering
\caption{Atribut Umum EV10 Redis Operation}
\label{tab:ev10-umum}
\renewcommand{\arraystretch}{1.3}
\begin{tabular}{|p{3.5cm}|p{8.5cm}|}
\hline
\textbf{Atribut} & \textbf{Kemungkinan Nilai} \\
\hline
\texttt{source\_component} &
\texttt{ProxyReencryption} \\
\hline
\texttt{actor\_id} &
\texttt{proxy\_iota\_address} PRE Server yang mengeksekusi operasi Redis. \\
\hline
\texttt{target\_object} &
Kunci Redis yang dioperasikan, misalnya
\texttt{nonce:<iota\_address>} atau
\texttt{keys:<hp\_addr>@<patient\_addr>}. \\
\hline
\texttt{action\_type} &
\texttt{CREATE} | \texttt{READ} | \texttt{DELETE} \\
\hline
\texttt{outcome} &
\texttt{Success} | \texttt{Failure} \\
\hline
\end{tabular}
\end{table}

\begin{table}[htbp]
\centering
\caption{Atribut Tambahan EV10 Redis Operation}
\label{tab:ev10-detail}
\renewcommand{\arraystretch}{1.3}
\begin{tabular}{|p{4cm}|p{7.5cm}|}
\hline
\textbf{Atribut} & \textbf{Kemungkinan Nilai dan Keterangan} \\
\hline
\texttt{redis\_key\_type} &
\texttt{nonce} | \texttt{access\_keys} --- kategori objek Redis yang
dioperasikan. \\
\hline
\texttt{operation\_type} &
\texttt{SET} | \texttt{GET} | \texttt{DEL} | \texttt{EXPIRE} ---
operasi Redis spesifik yang dieksekusi. \\
\hline
\texttt{ttl\_remaining} &
Sisa TTL entri Redis dalam detik pada saat operasi dilakukan;
\texttt{-1} apabila tidak ada TTL, \texttt{-2} apabila kunci tidak ditemukan. \\
\hline
\texttt{key\_pattern} &
Pola kunci Redis yang digunakan, misalnya \texttt{nonce:*} atau
\texttt{keys:*@*}, untuk membantu korelasi tanpa mengekspos nilai kunci
penuh. \\
\hline
\end{tabular}
\end{table}

% -----------------------------------------------------------------------
\subsubsection{EV11 --- IPFS Operation}
\label{subsec:ev11-ipfs-operation}

EV11 mencatat setiap operasi pengunggahan maupun pengunduhan objek pada IPFS
Cluster yang dilakukan oleh PRE Server. Objek yang dikelola adalah konten
klinis terenkripsi (AES-GCM) milik pasien. Pencatatan \textit{event} ini
diperlukan untuk mendeteksi pengunduhan data klinis oleh pihak tidak berwenang
maupun pembanjiran permintaan ke IPFS (T21, T22).

\begin{table}[htbp]
\centering
\caption{Atribut Umum EV11 IPFS Operation}
\label{tab:ev11-umum}
\renewcommand{\arraystretch}{1.3}
\begin{tabular}{|p{3.5cm}|p{8.5cm}|}
\hline
\textbf{Atribut} & \textbf{Kemungkinan Nilai} \\
\hline
\texttt{source\_component} &
\texttt{ProxyReencryption} \\
\hline
\texttt{actor\_id} &
\texttt{proxy\_iota\_address} PRE Server yang mengeksekusi operasi IPFS. \\
\hline
\texttt{target\_object} &
CID IPFS dari konten yang diunggah atau diunduh. \\
\hline
\texttt{action\_type} &
\texttt{CREATE} | \texttt{READ} \\
\hline
\texttt{outcome} &
\texttt{Success} | \texttt{Failure} \\
\hline
\end{tabular}
\end{table}

\begin{table}[htbp]
\centering
\caption{Atribut Tambahan EV11 IPFS Operation}
\label{tab:ev11-detail}
\renewcommand{\arraystretch}{1.3}
\begin{tabular}{|p{4cm}|p{7.5cm}|}
\hline
\textbf{Atribut} & \textbf{Kemungkinan Nilai dan Keterangan} \\
\hline
\texttt{cid} &
\textit{Content Identifier} IPFS dari objek yang dioperasikan. \\
\hline
\texttt{operation\_type} &
\texttt{upload} | \texttt{download} --- arah operasi terhadap IPFS. \\
\hline
\texttt{data\_size} &
Ukuran data dalam \textit{byte} yang diunggah atau diunduh. \\
\hline
\texttt{patient\_iota\_address} &
IOTA \textit{address} pasien pemilik data klinis yang dioperasikan. \\
\hline
\texttt{ipfs\_node\_url} &
URL \textit{node} IPFS Cluster yang digunakan dalam operasi. \\
\hline
\end{tabular}
\end{table}

% -----------------------------------------------------------------------
\subsubsection{EV12 --- IOTA Metadata Operation}
\label{subsec:ev12-iota-metadata}

EV12 mencatat setiap aktivitas pembacaan atau pembaruan objek \textit{on-chain}
pada jaringan IOTA, yang dapat dilakukan oleh seluruh komponen sistem. Objek
yang dicakup meliputi \texttt{AddressId}, \texttt{PatientIdAccount},
\texttt{HospitalPersonnelIdAccount}, \texttt{HospitalIdMetadata}, dan
objek-objek \textit{shared} lainnya. Pencatatan \textit{event} ini diperlukan
untuk mendeteksi potensi analisis pola transaksi dari metadata publik (T23).

\begin{table}[htbp]
\centering
\caption{Atribut Umum EV12 IOTA Metadata Operation}
\label{tab:ev12-umum}
\renewcommand{\arraystretch}{1.3}
\begin{tabular}{|p{3.5cm}|p{8.5cm}|}
\hline
\textbf{Atribut} & \textbf{Kemungkinan Nilai} \\
\hline
\texttt{source\_component} &
\texttt{HospitalTauriBackend} | \texttt{PatientTauriBackend} |
\texttt{MinistryTauriBackend} | \texttt{ProxyReencryption} \\
\hline
\texttt{actor\_id} &
IOTA \textit{address} pengguna | \texttt{proxy\_iota\_address} apabila
PRE Server. \\
\hline
\texttt{target\_object} &
IOTA Object ID dari objek \textit{on-chain} yang diakses. \\
\hline
\texttt{action\_type} &
\texttt{READ} | \texttt{UPDATE} \\
\hline
\texttt{outcome} &
\texttt{Success} | \texttt{Failure} \\
\hline
\end{tabular}
\end{table}

\begin{table}[htbp]
\centering
\caption{Atribut Tambahan EV12 IOTA Metadata Operation}
\label{tab:ev12-detail}
\renewcommand{\arraystretch}{1.3}
\begin{tabular}{|p{4cm}|p{7.5cm}|}
\hline
\textbf{Atribut} & \textbf{Kemungkinan Nilai dan Keterangan} \\
\hline
\texttt{object\_id} &
IOTA Object ID dari objek yang diakses atau diperbarui. \\
\hline
\texttt{object\_type} &
\texttt{AddressId} | \texttt{PatientIdAccount} |
\texttt{HospitalPersonnelIdAccount} | \texttt{HospitalIdMetadata} |
tipe objek lainnya. \\
\hline
\texttt{move\_function} &
Nama fungsi Move yang digunakan untuk mengakses objek, apabila akses
dilakukan melalui \texttt{dev\_inspect} atau transaksi. \\
\hline
\texttt{is\_mutable} &
\texttt{true} | \texttt{false} --- menyatakan apakah operasi bersifat
mutasi terhadap \textit{state on-chain} atau hanya pembacaan. \\
\hline
\texttt{transaction\_digest} &
\textit{Digest} transaksi IOTA apabila operasi bersifat mutasi;
\texttt{null} apabila hanya \texttt{dev\_inspect}. \\
\hline
\texttt{dev\_inspect\_used} &
\texttt{true} | \texttt{false} --- menyatakan apakah pembacaan dilakukan
menggunakan mekanisme \texttt{dev\_inspect} tanpa menerbitkan transaksi. \\
\hline
\end{tabular}
\end{table}

% -----------------------------------------------------------------------
\subsubsection{EV13 --- Capability Validation}
\label{subsec:ev13-capability-validation}

EV13 mencatat setiap proses validasi \textit{capability} yang dilakukan oleh
PRE Server sebelum permintaan sensitif diproses lebih lanjut. Objek yang
divalidasi meliputi \texttt{ProxyCapEnum} dan \texttt{GlobalAdminCap}
(\textit{on-chain}), serta JWT \textit{token} yang diterbitkan PRE Server untuk
\textit{hospital personnel}. Pencatatan \textit{event} ini diperlukan untuk
mendeteksi upaya eskalasi privilese melalui kelemahan validasi \textit{capability}
(T24, T25).

\begin{table}[htbp]
\centering
\caption{Atribut Umum EV13 Capability Validation}
\label{tab:ev13-umum}
\renewcommand{\arraystretch}{1.3}
\begin{tabular}{|p{3.5cm}|p{8.5cm}|}
\hline
\textbf{Atribut} & \textbf{Kemungkinan Nilai} \\
\hline
\texttt{source\_component} &
\texttt{ProxyReencryption} \\
\hline
\texttt{actor\_id} &
IOTA \textit{address} \textit{hospital personnel} yang divalidasi |
\texttt{proxy\_iota\_address} untuk validasi \texttt{ProxyCap}. \\
\hline
\texttt{target\_object} &
IOTA Object ID \textit{capability on-chain} | JWT \textit{token identifier}. \\
\hline
\texttt{action\_type} &
\texttt{VALIDATE} \\
\hline
\texttt{outcome} &
\texttt{Success} | \texttt{Denied} | \texttt{Failure} \\
\hline
\end{tabular}
\end{table}

\begin{table}[htbp]
\centering
\caption{Atribut Tambahan EV13 Capability Validation}
\label{tab:ev13-detail}
\renewcommand{\arraystretch}{1.3}
\begin{tabular}{|p{4cm}|p{7.5cm}|}
\hline
\textbf{Atribut} & \textbf{Kemungkinan Nilai dan Keterangan} \\
\hline
\texttt{capability\_id} &
IOTA Object ID \textit{capability on-chain} atau pengenal JWT yang
divalidasi. \\
\hline
\texttt{required\_scope} &
Cakupan akses yang dipersyaratkan oleh operasi yang diminta, misalnya
\texttt{read} atau \texttt{update}. \\
\hline
\texttt{actual\_scope} &
Cakupan akses aktual yang dimiliki oleh \textit{capability} yang
diperiksa. \\
\hline
\texttt{rejection\_reason} &
\texttt{expired} | \texttt{scope\_mismatch} | \texttt{not\_found} |
\texttt{invalid\_signature} | \texttt{null} apabila validasi berhasil. \\
\hline
\texttt{middleware\_layer} &
\texttt{jwt\_middleware} | \texttt{onchain\_check} --- lapisan validasi
yang melakukan pemeriksaan \textit{capability}. \\
\hline
\end{tabular}
\end{table}

\subsection{Rancangan Mekanisme Pembangkitan dan Penerimaan Audit Event}
\label{subsec:rancangan-pembangkitan-event}

Berdasarkan analisis pada Subbab~\ref{sec:analisis-event-source} dan
analisis mekanisme pengiriman log pada
Subbab~\ref{subsubsec:analisis-mekanisme-pengiriman-log}, \textit{audit
event} dibangkitkan dari empat \textit{event source} yang telah
diidentifikasi, yaitu \textit{Patient Client}, \textit{Hospital Client},
\textit{Ministry Client}, dan PRE Server. Rancangan mekanisme ini mencakup
tiga aspek utama yang saling berkaitan: enkapsulasi \textit{event} dalam
format terenkripsi dan tertandatangani untuk menjamin kerahasiaan dan
\textit{non-repudiation}, mekanisme pengiriman dengan jaminan ketercapaian
untuk memastikan tidak ada \textit{event} yang hilang, serta penerimaan dan
verifikasi pada sisi ATS.

\subsubsection{Rancangan Enkapsulasi Audit Event}
\label{subsubsec:rancangan-enkapsulasi-event}

Analisis pada Subbab~\ref{subsubsec:analisis-jaminan-keasilan-log}
menetapkan bahwa setiap \textit{event source} harus membungkus
\textit{audit event} dalam format yang menjamin keaslian pengirim dan
\textit{non-repu- diation}. Namun, terdapat kebutuhan keamanan tambahan yang
perlu dipenuhi, yaitu kerahasiaan isi \textit{event} selama transit. \textit{Audit event} memuat atribut
yang berpotensi sensitif, seperti identitas aktor, jenis tindakan yang
dilakukan, dan informasi objek yang diakses. Meskipun kanal komunikasi
umumnya dilindungi oleh TLS, perlindungan pada lapisan transport bersifat
\textit{hop-by-hop} dan tidak memberikan jaminan kerahasiaan yang melekat
pada data itu sendiri setelah transit selesai.

Untuk memenuhi kedua kebutuhan tersebut secara bersamaan, rancangan
mengadopsi pendekatan \textit{Encrypt-then-Sign}, yaitu \textit{payload
event} dienkripsi terlebih dahulu kemudian \textit{ciphertext} yang
dihasilkan ditandatangani. Pola ini dipilih karena tanda tangan diberikan
atas \textit{ciphertext}, bukan \textit{plaintext}, sehingga secara
kriptografis membuktikan bahwa pengirim yang menciptakan \textit{ciphertext}
tersebut dan bertanggung jawab atas isinya, bukan sekadar memiliki akses
atas informasi yang terkandung di dalamnya. Rancangan enkapsulasi terdiri
atas dua lapisan mekanisme yang diuraikan pada subbab-subbab berikut.

\paragraph{Lapisan Enkripsi: AES-256-GCM dengan Distribusi Kunci ECIES}

\textit{Payload audit event} yang telah diserialisasi ke JSON dienkripsi
menggunakan algoritma AES-256-GCM dengan kunci sesi yang dibangkitkan secara
acak pada setiap pengiriman. AES-256-GCM dipilih karena bersifat
\textit{authenticated encryption}, yaitu selain menjamin kerahasiaan,
algoritma ini juga menghasilkan \textit{authentication tag} yang memastikan
integritas \textit{ciphertext} tidak dapat dirusak tanpa terdeteksi.

Distribusi kunci sesi AES ke ATS dilakukan menggunakan mekanisme ECIES
(\textit{Elliptic Curve Integrated Encryption Scheme}) berbasis kurva P-256.
Alur distribusi kunci dirancang sebagai berikut:

\begin{enumerate}
    \item Pengirim membangkitkan pasangan kunci \textit{ephemeral} P-256
    untuk setiap pengiriman.
    \item Pengirim melakukan operasi ECDH antara kunci privat \textit{ephemeral}
    dengan kunci publik ATS yang telah dikonfigurasi sebelumnya, menghasilkan
    \textit{shared secret}.
    \item \textit{Shared secret} diturunkan menjadi kunci pembungkus
    (\textit{wrapping key}) 256-bit menggunakan KDF SHA-256.
    \item Kunci sesi AES dienkripsi menggunakan kunci pembungkus tersebut
    dengan AES-256-GCM, menghasilkan kunci sesi terenkripsi.
    \item Kunci publik \textit{ephemeral}, \textit{nonce} pembungkus, dan
    kunci sesi terenkripsi digabungkan menjadi satu bidang \texttt{enc\_aes\_key}
    dalam format \texttt{<base64(epheme- ral\_pubkey)>.<base64(wrap\_nonce)>
    .<base64(enc\_aes\_key)>}.
\end{enumerate}

Dengan skema ini, kunci sesi AES hanya dapat dipulihkan oleh ATS menggunakan
kunci privatnya melalui operasi ECDH yang berkebalikan, tanpa memerlukan
pertukaran kunci awal (\textit{key exchange}) yang terpisah antara pengirim
dan penerima. Setiap pengiriman menggunakan kunci \textit{ephemeral} yang
berbeda, sehingga kompromi atas satu kunci sesi tidak mengancam kerahasiaan
\textit{event} lainnya (\textit{forward secrecy}).

\paragraph{Lapisan Penandatanganan: Ed25519 dengan Kunci Privat IOTA Aktor}

Setelah \textit{payload} dienkripsi, \textit{ciphertext} yang dihasilkan
ditandatangani menggunakan kunci privat IOTA Ed25519 dari aktor yang memicu
\textit{event}. Penggunaan kunci privat IOTA aktor, bukan kunci yang
dibangkitkan sementara oleh modul ATS pada komponen pengirim, merupakan
keputusan rancangan yang krusial untuk memperkuat sifat
\textit{non-repudiation}.

Dasar pertimbangannya adalah sebagai berikut. Kunci privat IOTA setiap aktor
telah terikat secara kriptografis dengan identitas aktor melalui mekanisme
BIP39 dan derivasi kunci yang digunakan DecMed; setiap aktor memiliki satu
pasang kunci yang konsisten sepanjang siklus hidupnya dalam sistem dan yang
sama digunakannya untuk menandatangani transaksi IOTA yang tercatat secara
\textit{on-chain}. Dengan demikian, tanda tangan pada \textit{audit event}
dapat dihubungkan langsung ke \textit{IOTA address} pengirim yang tercatat
pada \textit{ledger}. Apabila suatu aktor menyangkal telah membangkitkan
sebuah \textit{audit event}, tanda tangan yang ada pada \textit{event}
tersebut dapat diverifikasi terhadap kunci publik yang terhubung ke
\textit{IOTA address} yang sama yang digunakan aktor tersebut untuk
bertransaksi di jaringan IOTA, sehingga penyangkalan menjadi tidak dapat
diterima selama kunci privat tidak dikompromikan.

Tanda tangan Ed25519 dibuat menggunakan skema \textit{personal message} yang
konsisten dengan mekanisme penandatanganan pada ekosistem IOTA, yaitu pesan
yang ditandatangani adalah \textit{byte-byte ciphertext} yang dibungkus dalam
\texttt{IntentMessage} dengan \texttt{Intent::personal\_message()}. Hasilnya
diserialisasi ke dalam representasi Base64.

\paragraph{Format Akhir: EncryptedSignedEvent}

Hasil akhir enkapsulasi adalah objek \texttt{EncryptedSignedEvent} yang
memuat lima bidang sebagaimana ditunjukkan pada
Tabel~\ref{tab:encrypted-signed-event}.

\begin{table}[htbp]
\centering
\caption{Bidang pada Struktur \texttt{EncryptedSignedEvent}}
\label{tab:encrypted-signed-event}
\renewcommand{\arraystretch}{1.3}
\begin{tabular}{|p{3.5cm}|p{8.5cm}|}
\hline
\textbf{Bidang} & \textbf{Deskripsi} \\
\hline
\texttt{enc\_aes\_key} &
Kunci sesi AES yang dienkripsi dengan ECIES, dalam format
\texttt{<ephemeral\_pubkey>.<wrap\_nonce>.<enc\_key>} yang masing-masing
dikodekan Base64 dan dipisahkan titik. \\
\hline
\texttt{ciphertext} &
\textit{Payload audit event} yang dienkripsi menggunakan AES-256-GCM,
dikodekan Base64. \\
\hline
\texttt{nonce} &
\textit{Nonce} AES-GCM sepanjang 12 byte, dikodekan Base64. \\
\hline
\texttt{signature} &
Tanda tangan Ed25519 atas \textit{byte-byte ciphertext}, dikodekan Base64. \\
\hline
\texttt{iota\_address} &
\textit{IOTA address} pengirim yang terhubung ke kunci publik Ed25519 yang
digunakan untuk membuat tanda tangan, digunakan oleh ATS untuk keperluan
verifikasi. \\
\hline
\end{tabular}
\end{table}

Skema ini memisahkan secara tegas antara mekanisme kerahasiaan (ECIES +
AES-256-GCM) dan mekanisme \textit{non-repudiation} (Ed25519 atas
\textit{ciphertext}), sehingga kedua properti keamanan dapat diverifikasi
secara independen oleh pihak yang berwenang.

\subsubsection{Rancangan Mekanisme Pengiriman dengan Jaminan Ketercapaian}
\label{subsubsec:rancangan-pengiriman-event}

Analisis pada Subbab~\ref{subsubsec:analisis-mekanisme-pengiriman-log}
menetapkan bahwa pengiriman \textit{event} harus bersifat
\textit{asynchronous push-based} agar tidak mengganggu latensi layanan
utama. Namun, sifat asinkron menimbulkan risiko kehilangan \textit{event}
apabila ATS tidak dapat dijangkau pada saat \textit{event} dibangkitkan,
misalnya karena gangguan jaringan sementara atau ATS sedang tidak beroperasi.
Kegagalan pencatatan \textit{event} yang tidak terdeteksi merupakan ancaman
terhadap kelengkapan audit trail sebagaimana disyaratkan oleh KNF-TA-04,
sehingga mekanisme pengiriman harus memberikan jaminan ketercapaian
(\textit{delivery guarantee}) yang kuat di samping tetap bersifat
\textit{non-blocking}.

Untuk menjawab kebutuhan ini, rancangan mengadopsi pola
\textit{persistent queue} pada sisi pengirim yang digabungkan dengan
pengiriman asinkron. Perbandingan antara pendekatan pengiriman murni
\textit{fire-and-forget} dan pendekatan \textit{persistent queue} yang
diusulkan ditunjukkan pada Tabel~\ref{tab:analisis-pengiriman-event}.

\begin{table}[htbp]
\centering
\caption{Analisis Pendekatan Mekanisme Pengiriman \textit{Audit Event}}
\label{tab:analisis-pengiriman-event}
\renewcommand{\arraystretch}{1.3}
\begin{tabular}{|p{3cm}|p{4.5cm}|p{5cm}|}
\hline
\textbf{Pendekatan} & \textbf{Keunggulan} & \textbf{Keterbatasan} \\
\hline
\textit{Fire-and-Forget} Murni &
\textit{Non-blocking}, tidak menambah kompleksitas komponen pengirim. &
\textit{Event} hilang secara permanen apabila ATS tidak dapat dijangkau
saat pengiriman; tidak ada mekanisme pemulihan. \\
\hline
\textit{Persistent Queue} + Pengiriman Asinkron &
Menjamin setiap \textit{event} akan terkirim pada akhirnya melalui
mekanisme \textit{retry}; tahan terhadap gangguan jaringan sementara dan
\textit{restart} komponen. &
Membutuhkan pengelolaan berkas antrian lokal dan \textit{background worker}
pada setiap \textit{event source}. \\
\hline
\end{tabular}
\end{table}

Berdasarkan Tabel~\ref{tab:analisis-pengiriman-event}, pendekatan
\textit{persistent queue} dipilih karena jaminan ketercapaian yang
diberikannya jauh lebih penting bagi integritas audit trail dibandingkan
kompleksitas implementasi tambahan yang ditimbulkannya. Alur pengiriman
\textit{event} yang dirancang adalah sebagai berikut.

\begin{enumerate}
    \item \textbf{Konstruksi \texttt{EncryptedSignedEvent}.} Setiap kali
    sebuah \textit{event} dibangkitkan, komponen modul \texttt{ats} pada
    \textit{event source} mengonstruksi \texttt{EncryptedSignedEvent} sesuai
    mekanisme enkapsulasi pada
    Subbab~\ref{subsubsec:rancangan-enkapsulasi-event}.

    \item \textbf{Penyimpanan ke antrian persisten.} Sebelum percobaan
    pengiriman apa pun dilakukan, \texttt{EncryptedSignedEvent} yang telah
    dikonstruksi disimpan terlebih dahulu ke dalam berkas antrian persisten
    berbasis \textit{append-only} pada penyimpanan lokal \textit{event source}.
    Berkas antrian menggunakan format JSON-Lines, di mana setiap baris merupakan
    satu entri antrian (\texttt{QueueEntry}) yang memuat \textit{signed payload},
    pengenal unik entri, label konteks, waktu pembuatan, dan jumlah percobaan
    pengiriman yang telah dilakukan. Penyimpanan ke antrian dilakukan terlebih
    dahulu agar \textit{event} aman tersimpan secara lokal terlepas dari hasil
    pengiriman berikutnya.

    \item \textbf{Percobaan pengiriman langsung.} Setelah entri berhasil
    tersimpan di antrian, pengirim segera mencoba mengirimkan
    \texttt{EncryptedSignedEvent} ke \textit{endpoint}
    \texttt{POST /api/events} pada ATS secara langsung dalam \textit{task}
    asinkron yang tidak memblokir alur layanan utama. Apabila pengiriman
    langsung berhasil (ATS mengembalikan respons sukses), entri yang
    bersangkutan dihapus dari antrian.

    \item \textbf{Pengiriman ulang oleh \textit{background retry worker}.}
    Apabila pengiriman langsung gagal karena ATS tidak dapat dijangkau atau
    mengembalikan respons \textit{error}, entri tetap berada di antrian.
    \textit{Background retry worker} yang berjalan secara periodik pada setiap
    \textit{event source} membaca seluruh entri yang tersisa di antrian dan
    mencoba mengirimkannya kembali. \textit{Worker} ini menggunakan strategi
    \textit{exponential backoff} untuk menentukan interval antar-percobaan,
    sehingga menghindari pembanjiran ATS saat baru pulih dari gangguan.
    Entri yang telah melebihi batas percobaan pengiriman dikecualikan
    sementara dari percobaan ulang, tetapi tetap disimpan di antrian agar
    tidak hilang secara permanen dan dapat diinspeksi oleh operator.
\end{enumerate}

Rancangan ini menjamin bahwa setiap \textit{event} yang berhasil dikonstruksi
dan disimpan ke antrian akan terkirim ke ATS pada akhirnya, bahkan jika
\textit{event source} mengalami \textit{restart} di antara percobaan
pengiriman, selama berkas antrian tidak rusak secara fisik. Pemisahan antara
tahap penyimpanan ke antrian dan tahap pengiriman ke ATS juga memastikan bahwa
kegagalan jaringan tidak menyebabkan \textit{event} hilang secara diam-diam
tanpa sepengetahuan operator.

\subsubsection{Rancangan Penerimaan dan Verifikasi Event pada ATS}
\label{subsubsec:rancangan-penerimaan-event}

Pada sisi ATS, \texttt{EncryptedSignedEvent} yang diterima melalui
\textit{endpoint} \texttt{POST /api/events} diproses dalam dua fase sebelum
diteruskan ke komponen Audit Logger.

\paragraph{Fase Pertama: Verifikasi Tanda Tangan}

ATS mengekstrak \texttt{iota\_address} dari \textit{event} yang diterima,
kemudian memverifikasi tanda tangan Ed25519 yang terdapat pada bidang
\texttt{signature} terhadap \textit{byte-byte ciphertext} pada bidang
\texttt{ciphertext}. Verifikasi menggunakan kunci publik Ed25519 yang
diturunkan dari \texttt{iota\_address} pengirim. Apabila verifikasi
berhasil, ATS dapat menyatakan bahwa \textit{event} berasal dari pemegang
kunci privat yang terhubung ke \texttt{iota\_address} yang diklaim, dan bahwa
\textit{ciphertext} tidak mengalami perubahan sejak ditandatangani. Apabila
verifikasi gagal, \textit{event} ditolak tanpa diproses lebih lanjut,
mencegah injeksi \textit{event} palsu ke dalam rantai audit.

Pada fase ini, ATS tidak mendekripsi isi \textit{event}. Verifikasi tanda
tangan dilakukan sepenuhnya terhadap \textit{ciphertext}, konsisten dengan
pola \textit{Encrypt-then-Sign} yang diadopsi. Keputusan untuk memisahkan
fase verifikasi dari fase dekripsi merupakan keputusan rancangan yang disengaja:
verifikasi adalah hak ATS sebagai penjaga gerbang pencatatan, sedangkan
dekripsi adalah hak pihak yang berwenang melakukan audit atau investigasi.

\paragraph{Fase Kedua: Penerusan ke Antrian Internal}

\texttt{EncryptedSignedEvent} yang telah terverifikasi tanda tangannya
dikirimkan ke antrian asinkron internal yang menghubungkan komponen Audit
Receiver dengan komponen Audit Logger. Antrian ini diimplementasikan sebagai
kanal \textit{multi-producer single-consumer} (mpsc) dengan kapasitas yang
cukup untuk menyerap lonjakan kedatangan \textit{event}. Desain ini memisahkan
proses penerimaan HTTP dari proses pembentukan dan penulisan \textit{audit
record}, sehingga keduanya dapat berjalan secara konkuren tanpa saling
memblokir. Antrian internal bertindak sebagai penyangga (\textit{buffer})
yang mencegah lonjakan \textit{event} menghambat penerimaan \textit{event}
baru berikutnya.

\subsubsection{Rancangan Penyimpanan Audit Record dalam Kondisi Terenkripsi}
\label{subsubsec:rancangan-penyimpanan-terenkripsi}

Setelah \texttt{EncryptedSignedEvent} terverifikasi dan diteruskan oleh Audit
Logger, \textit{event} tersebut disimpan ke dalam \textit{audit record} dan
kemudian ke berkas log lokal serta IPFS dalam kondisi masih terenkripsi.
Keputusan rancangan untuk tidak mendekripsi \textit{event} pada titik ini
didasarkan pada dua pertimbangan berikut.

Pertama, pemisahan tanggung jawab. ATS berperan sebagai penyimpan bukti
digital yang andal, bukan sebagai pemroses isi \textit{event}. Mendekripsi
\textit{event} sebelum penyimpanan akan memperluas permukaan serangan karena
kunci privat ATS harus tersedia secara aktif setiap kali \textit{event}
diproses. Dengan menyimpan \textit{event} dalam kondisi terenkripsi, kunci
privat ATS hanya perlu tersedia saat audit atau investigasi dilakukan.

Kedua, kerahasiaan jangka panjang. Berkas log yang tersimpan di IPFS dapat
diakses oleh pihak mana pun yang mengetahui CID-nya. Apabila \textit{event}
disimpan dalam kondisi \textit{plaintext}, kerahasiaan isi audit trail
bergantung sepenuhnya pada keamanan akses ke IPFS. Dengan menyimpan
\textit{event} dalam kondisi terenkripsi, kerahasiaan isi audit trail
dijamin secara kriptografis: hanya pemegang kunci privat ATS yang dapat
mendekripsi dan membaca isi \textit{event}, terlepas dari siapa yang dapat
mengakses berkas log di IPFS.

Implikasi praktis dari keputusan ini adalah bahwa mekanisme
\textit{hash-chaining} pada Subbab~\ref{subsec:rancangan-pembentukan-record}
diterapkan terhadap \texttt{EncryptedSignedEvent} sebagai satu kesatuan,
bukan terhadap isi \textit{plaintext event}. Nilai \texttt{record\_hash}
dihitung dari kombinasi \texttt{record\_id}, \texttt{timestamp},
\texttt{prev\_record\_hash}, dan seluruh bidang \texttt{EncryptedSignedEvent}
yang mencakup \texttt{ciphertext} dan \texttt{signature}. Dengan demikian,
integritas rantai \textit{audit record} tetap dapat diverifikasi tanpa perlu
mendekripsi isi \textit{event} terlebih dahulu: perubahan pada
\textit{ciphertext} mana pun akan merusak \texttt{record\_hash} dan
mengakibatkan kegagalan verifikasi rantai, sebagaimana berlaku pada
mekanisme \textit{hash-chaining} yang telah dianalisis pada
Subbab~\ref{subsubsec:analisis-hash-chaining}.

\subsection{Rancangan Mekanisme Pembentukan dan Penyimpanan Audit Record}
\label{subsec:rancangan-pembentukan-record}

Setelah \textit{event} yang telah terverifikasi diteruskan dari antrian internal, komponen Audit Logger membentuk sebuah \texttt{AuditRecord} yang terdiri atas komponen-komponen berikut.

\begin{itemize}
    \item \texttt{record\_id}, berupa pengenal unik berbasis UUID v7 yang mengandung informasi waktu sehingga bersifat monoton dan dapat diurutkan secara kronologis.
    \item \texttt{timestamp}, berupa waktu pembentukan \textit{record} dalam format UTC yang diambil pada saat Audit Logger memproses \textit{event}.
    \item \texttt{prev\_record\_hash}, berupa nilai \textit{hash} SHA-256 dari \textit{record} sebelumnya, atau \texttt{null} apabila merupakan \textit{record} pertama dalam rantai.
    \item \texttt{record\_hash}, berupa nilai \textit{hash} SHA-256 yang dihitung dari kombinasi \texttt{record\_id}, \texttt{timestamp}, \texttt{prev\_record\_hash}, dan isi \textit{event}.
    \item Seluruh atribut \texttt{AuditEvent} yang telah terverifikasi.
\end{itemize}

Mekanisme \textit{hash-chaining} pada \textit{audit record} memberikan sifat
\textit{tamper-evident} pada berkas log lokal sebelum berkas tersebut diarsipkan
ke IPFS. Nilai \texttt{record\_hash} yang menyertakan \texttt{prev\_record\_hash}
menjadikan setiap \textit{record} terikat secara kriptografis pada seluruh
\textit{record} sebelumnya dalam rantai, sebagaimana dianalisis pada
Subbab~\ref{subsubsec:analisis-hash-chaining}. Modifikasi \textit{record},
penyisipan \textit{record}, maupun penghapusan \textit{record} yang masih diikuti
oleh \textit{record} lain akan menyebabkan ketidaksesuaian antara
\texttt{prev\_record\_hash} dan \texttt{record\_hash} pada \textit{record}
sesudahnya, sehingga perusakan dapat terdeteksi melalui pemeriksaan konsistensi
rantai.
 
Namun, \textit{hash chain} hanya dapat membuktikan konsistensi internal antar
\textit{record} yang masih ada. Penghapusan satu atau beberapa \textit{record}
terakhir pada berkas log aktif menghasilkan rantai yang lebih pendek tetapi
tetap konsisten, karena tidak ada \textit{record} sesudahnya yang merujuk pada
\textit{record} yang dihapus. Penghapusan semacam ini hanya dapat terdeteksi
apabila nilai \textit{hash} ujung rantai telah dicatat pada media di luar
kendali ALS. Pada rancangan ini, peran tersebut dijalankan oleh
\texttt{final\_record\_hash} dan \texttt{record\_count} pada \texttt{AuditBatch}
yang dipublikasikan sebagai \textit{frozen object} di IOTA pada saat rotasi,
sebagaimana diuraikan pada Subbab~\ref{subsec:rancangan-rotasi-log}. Dengan
demikian, \textit{record} yang telah melalui rotasi terlindungi dari
penghapusan pada bagian ujung, sedangkan \textit{record} pada berkas log aktif
yang belum dirotasi hanya terlindungi dari modifikasi, penyisipan, dan
penghapusan di tengah rantai.

Setiap \texttt{AuditRecord} yang terbentuk kemudian ditulis secara \textit{append} ke berkas log lokal dalam format JSON-Lines (satu objek JSON per baris), menjamin bahwa berkas log tidak pernah menimpa entri yang sudah ada. Hal ini merealisasikan sifat \textit{append-only} yang disyaratkan oleh kebutuhan nonfungsional KNF-TA-01 pada Tabel~\ref{tab:kebutuhan-nonfungsional}.

\subsection{Rancangan Mekanisme Rotasi Berkas Log}
\label{subsec:rancangan-rotasi-log}

Berdasarkan analisis mekanisme penyimpanan berkas log pada Subbab~\ref{subsubsec:analisis-penyimpanan-berkas}, IPFS menggunakan pengalamatan berbasis konten (\textit{content-addressed storage}) sehingga setiap berkas diidentifikasi oleh \textit{hash} dari isinya. Konsekuensinya, \textit{audit record} tidak dapat diunggah satu per satu ke IPFS secara langsung, karena setiap unggahan terpisah menghasilkan CID yang berbeda-beda dan tidak membentuk satu arsip yang dapat diverifikasi secara utuh. Oleh karena itu, diperlukan mekanisme pengelompokan dan perputaran berkas log (\textit{log rotation}) yang mengelola berkas log lokal sebagai \textit{buffer} sementara sebelum pengarsipan ke IPFS.

Komponen Log Rotation \& Archival Worker dirancang untuk berjalan sebagai \textit{background task} yang berulang secara periodik pada interval yang dapat dikonfigurasi melalui konstanta \texttt{LOG\_ROTATION\_INTERVAL\_SECS}. Pada setiap siklus rotasi, worker terlebih dahulu memeriksa apakah berkas log aktif ada dan tidak kosong. Apabila berkas kosong, siklus dilewati untuk menghindari pencatatan metadata rotasi yang tidak bermakna pada jaringan IOTA.

Apabila berkas log aktif berisi data, worker membentuk nama berkas arsip baru dengan format \texttt{logs/uploading\_\{unix\_timestamp\}.log} menggunakan \textit{Unix timestamp} pada saat rotasi dilakukan, kemudian memindahkan berkas log aktif ke nama arsip baru. Pemindahan berkas dilakukan menggunakan operasi \texttt{rename} yang bersifat atomik pada tingkat sistem operasi, sehingga tidak ada jeda di mana berkas log aktif tidak tersedia untuk penulisan. Setelah pemindahan selesai, ALS langsung membuka berkas log aktif baru untuk menampung \textit{audit record} yang masuk selama proses pengarsipan berkas lama berlangsung.

Berkas dikirimkan menggunakan mekanisme \textit{streaming} berbasis \texttt{BytesCodec} sehingga isi berkas tidak perlu dimuat seluruhnya ke memori sebelum dikirim; hal ini penting untuk mengakomodasi berkas log yang berukuran besar tanpa membebani konsumsi memori ALS. IPFS Cluster mengembalikan respons yang memuat CID dari berkas yang berhasil diunggah. CID ini bersifat unik terhadap isi berkas dan menjadi penanda permanen yang digunakan pada tahap publikasi metadata ke IOTA. Apabila unggah gagal, siklus rotasi dihentikan pada tahap ini dan kesalahan dicatat ke \textit{stderr}; berkas arsip tetap berada di penyimpanan lokal dengan awalan \texttt{uploading\_} sehingga dapat diidentifikasi dan diproses ulang pada siklus rotasi berikutnya.

Berdasarkan analisis strategi persistensi data pada Subbab~\ref{subsubsec:analisis-persistensi-data}, berkas yang diunggah ke IPFS rentan terhapus oleh mekanisme \textit{Garbage Collection} apabila tidak ada penanda eksplisit yang menyatakan bahwa berkas tersebut perlu dipertahankan. Oleh karena itu, setiap berkas arsip log yang berhasil diunggah diterapkan \textit{pinning} pada IPFS Cluster melalui \textit{endpoint} \texttt{/pins} yang tersedia. \textit{Pinning} memberi penanda eksplisit kepada \textit{node} IPFS agar berkas log tersebut tidak diidentifikasi sebagai data tidak terpakai dan tidak terhapus saat proses \textit{Garbage Collection} berjalan. Retensi jangka panjang berkas arsip log bergantung pada kebijakan \textit{pinning} IPFS Cluster yang sama yang telah digunakan DecMed untuk data klinis pasien, sehingga rancangan ini tidak memperkenalkan infrastruktur persistensi baru.

\subsection{Rancangan Mekanisme Publikasi Metadata \textit{On-Chain} ke IOTA}
\label{subsec:rancangan-publikasi-iota}

Setelah berkas arsip log berhasil diunggah ke IPFS dan CID diperoleh, tahap terakhir adalah menerbitkan metadata rotasi log ke jaringan IOTA sebagai \textit{immutable trust anchor}. Rancangan mekanisme ini mencakup tiga aspek: skema metadata, mekanisme penerbitan melalui \textit{smart contract} Move, dan pembentukan rantai metadata \textit{on-chain}.

\subsubsection{Skema Metadata Rotasi Log}
\label{subsubsec:rancangan-skema-metadata}

Objek metadata rotasi log dirancang untuk memuat seluruh informasi yang diperlukan agar verifikasi keutuhan berkas arsip log dapat dilakukan secara independen di masa mendatang. Rancangan skema metadata ditunjukkan pada Tabel~\ref{tab:skema-metadata-iota}.

\begin{table}[htbp]
\centering
\caption{Skema Metadata Rotasi Log pada \textit{Smart Contract} IOTA}
\label{tab:skema-metadata-iota}
\renewcommand{\arraystretch}{1.3}
\begin{tabular}{|p{4cm}|p{3cm}|p{5.5cm}|}
\hline
\textbf{Bidang} & \textbf{Tipe Data} & \textbf{Deskripsi dan Fungsi} \\
\hline
\texttt{cid} & \texttt{String} & \textit{Content Identifier} IPFS dari berkas arsip log yang diunggah. Digunakan untuk mengambil kembali berkas arsip dari IPFS saat verifikasi. \\
\hline
\texttt{file\_hash} & \texttt{String} (hex) & Nilai \textit{hash} SHA-256 dari berkas arsip, dihitung secara lokal sebelum unggah. Digunakan sebagai referensi integritas independen dari IPFS untuk mendeteksi apabila isi berkas yang diambil tidak sesuai. \\
\hline
\texttt{log\_sequence\_number} & \texttt{u64} & Nomor urut rotasi yang bersifat monoton dan bertambah satu setiap siklus rotasi. Digunakan untuk menentukan urutan kronologis arsip log tanpa bergantung pada \textit{timestamp} semata. \\
\hline
\texttt{rotation\_time} & \texttt{String} & Waktu pelaksanaan rotasi dalam format UTC. Digunakan untuk merekonstruksi garis waktu pengarsipan log. \\
\hline
\texttt{prev\_tx\_digest} & \texttt{String} (opsional) & \textit{Transaction digest} publikasi metadata rotasi sebelumnya. Membentuk rantai antar entri metadata \textit{on-chain} yang analog dengan mekanisme \textit{hash-chaining} pada level \textit{audit record}. Bernilai \texttt{null} pada rotasi pertama. \\
\hline
\end{tabular}
\end{table}

Kombinasi \texttt{cid} dan \texttt{file\_hash} pada skema metadata memberikan dua lapis jaminan verifikasi yang saling melengkapi. Apabila \texttt{cid} digunakan untuk mengambil berkas dari IPFS, \texttt{file\_hash} digunakan untuk memverifikasi bahwa berkas yang diperoleh identik dengan berkas yang tercatat pada metadata. Apabila keduanya cocok, keutuhan berkas arsip log tersebut dapat dinyatakan terjaga sejak saat rotasi dilakukan.

\subsubsection{Penerbitan Metadata melalui \textit{Smart Contract} Move}
\label{subsubsec:rancangan-sc-move}

Metadata rotasi log diterbitkan ke jaringan IOTA melalui \textit{smart contract} Move yang didefinisikan pada modul \texttt{als::audit\_log}. Modul ini menyediakan fungsi \texttt{create\_log} yang menerima representasi JSON dari metadata rotasi, membuat objek \texttt{LogRecord} pada \textit{ledger} IOTA, kemudian langsung membekukannya melalui \texttt{transfer::freeze\_object}. Pembekuan objek menjadikan \texttt{LogRecord} bersifat \textit{immutable} secara penuh setelah transaksi final: tidak ada entitas mana pun, termasuk ALS itu sendiri, yang dapat mengubah isi \texttt{LogRecord} setelah diterbitkan. Rancangan ini secara langsung mengatasi masalah M-TA-03 yang mengidentifikasi bahwa mekanisme pencatatan eksisting pada DecMed masih menyediakan fungsi pembaruan pada \textit{smart contract} Move yang ada. Penerbitan dilakukan melalui Gas Station Server yang sudah tersedia pada DecMed, sehingga ALS tidak perlu mengelola aset IOTA secara mandiri. Apabila transaksi berhasil dieksekusi, worker memperoleh \textit{transaction digest} dan \textit{object ID} dari \texttt{LogRecord} yang baru dibuat.

Bidang \texttt{prev\_tx\_digest} pada skema metadata dirancang untuk membentuk rantai antar entri metadata \textit{on-chain} yang analog dengan mekanisme \textit{hash-chaining} pada level \textit{audit record} (Subbab~\ref{subsec:rancangan-pembentukan-record}). Setiap \texttt{LogRecord} yang diterbitkan ke IOTA membawa referensi eksplisit terhadap \texttt{LogRecord} sebelumnya, sehingga membentuk urutan yang dapat ditelusuri dari entri terbaru ke entri paling awal. Melalui penelusuran rantai ini, pihak yang berwenang dapat memverifikasi kelengkapan arsip log secara kronologis dengan mencocokkan \texttt{log\_sequence\_number} untuk memastikan tidak ada rotasi yang terlewat. Setelah publikasi metadata ke IOTA berhasil, worker memperbarui nilai \texttt{prev\_tx\_digest} yang tersimpan dalam memori untuk digunakan pada siklus rotasi berikutnya, kemudian menghapus berkas arsip dari penyimpanan lokal. Penghapusan berkas lokal ini dilakukan karena bukti integritas jangka panjang sudah tersedia melalui dua lapis mekanisme: berkas arsip tersimpan secara permanen di IPFS melalui \textit{pinning}, dan metadata dengan CID serta \texttt{file\_hash} tersimpan secara \textit{immutable} di IOTA.

\subsubsection{Rantai Metadata \textit{On-Chain}}

Properti utama yang dihasilkan oleh rantai metadata \textit{on-chain} adalah kemampuan untuk memverifikasi kelengkapan arsip log secara kronologis. Pihak yang berwenang dapat menelusuri seluruh rantai \texttt{LogRecord} dari yang terbaru mundur ke rotasi pertama menggunakan \texttt{prev\_tx\_digest}, kemudian mencocokkan \texttt{log\_sequence\_number} untuk memastikan tidak ada rotasi yang terlewat dalam rangkaian arsip. Skema rantai metadata \textit{on-chain} ditunjukkan pada Gambar~\ref{fig:rantai-metadata-onchain}.

\begin{figure}[htbp]
\centering
% \includegraphics[width=\textwidth]{chapter-3/onchain_metadata_chain.png}
\caption{Skema Rantai Metadata Rotasi Log \textit{On-Chain} pada IOTA}
\label{fig:rantai-metadata-onchain}
\end{figure}

Setelah publikasi metadata ke IOTA berhasil, worker memperbarui nilai \texttt{prev\_tx\_digest} yang tersimpan dalam memori untuk digunakan pada siklus rotasi berikutnya, kemudian menghapus berkas arsip dari penyimpanan lokal menggunakan \texttt{fs::remove\_file}. Penghapusan berkas lokal ini dilakukan karena bukti integritas jangka panjang sudah tersedia melalui dua lapis mekanisme: berkas arsip tersimpan secara permanen di IPFS melalui \textit{pinning}, dan metadata dengan CID serta \textit{file\_hash} tersimpan secara \textit{immutable} di IOTA. Penyimpanan berkas arsip secara lokal dalam jangka panjang tidak diperlukan dan justru akan menghabiskan kapasitas penyimpanan lokal ALS.

\subsection{Pemetaan Rancangan terhadap Solusi}
\label{subsec:pemetaan-rancangan-solusi}

Setelah rancangan solusi diuraikan, setiap rancangan dipetakan kembali terhadap solusi yang telah dianalisis pada Subbab~\ref{sec:analisis-solusi}. Melalui pemetaan ini, keterkaitan langsung setiap rancangan R-TA dengan solusi S-TA yang diusulkan ditunjukkan. Hasil pemetaan rancangan terhadap solusi ditunjukkan pada Tabel~\ref{tab:pemetaan-rancangan-solusi}.

\begin{table}[htbp]
\centering
\caption{Pemetaan Rancangan terhadap Solusi}
\label{tab:pemetaan-rancangan-solusi}
\renewcommand{\arraystretch}{1.3}
\begin{tabular}{|p{1.5cm}|p{8.5cm}|p{2.5cm}|}
\hline
\textbf{ID} & \textbf{Nama Rancangan} & \textbf{Mereali- sasikan Solusi} \\
\hline
R-TA-01 & Rancangan arsitektur solusi dengan penambahan ALS sebagai \textit{centralized audit collector} yang terintegrasi dengan infrastruktur IPFS dan IOTA yang sudah ada pada DecMed. & S-TA-01 \\
\hline
R-TA-02 & Rancangan skema \textit{audit record} bertingkat yang memisahkan atribut umum dan atribut tambahan spesifik per jenis \textit{event}. & S-TA-04 \\
\hline
R-TA-03 & Rancangan mekanisme pembangkitan dan penerimaan \textit{audit event} berbasis \textit{signed payload} Ed25519 dengan pengiriman \textit{asynchronous push-based}. & S-TA-01 \\
\hline
R-TA-04 & Rancangan mekanisme pembentukan \textit{audit record} dengan \textit{hash-chaining} antar \textit{record} untuk menjamin sifat \textit{tamper-evident} pada berkas log lokal. & S-TA-03 \\
\hline
R-TA-05 & Rancangan mekanisme rotasi log, pengarsipan berkas ke IPFS dengan \textit{pinning}, dan publikasi metadata \textit{immutable} ke IOTA melalui \textit{smart contract} Move. & S-TA-03, S-TA-05 \\
\hline
\end{tabular}
\end{table}

Tabel~\ref{tab:pemetaan-rancangan-solusi} memastikan setiap solusi S-TA memiliki rancangan R-TA yang merealisasikannya, sehingga ketertelusuran dari analisis solusi ke spesifikasi rancangan tetap terjaga. Selanjutnya, pemenuhan kebutuhan sistem terhadap rancangan yang telah ditetapkan ditunjukkan pada Tabel~\ref{tab:pemenuhan-kebutuhan}.

\begin{xltabular}{\textwidth}{|p{1.8cm}|X|X|}
\caption{Pemetaan Kebutuhan ke Mekanisme Rancangan ALS}
\label{tab:pemenuhan-kebutuhan}\\
\hline
\textbf{ID} & \textbf{Kebutuhan} & \textbf{Mekanisme pada Rancangan} \\
\hline
\endfirsthead
\hline
\textbf{ID} & \textbf{Kebutuhan} & \textbf{Mekanisme pada Rancangan} \\
\hline
\endhead
\hline
\endfoot

KF-TA-01 & Mengagregasi data log dari berbagai komponen ke satu repositori terpusat. & \textit{Endpoint} \texttt{POST /api/events} pada ALS sebagai titik masuk tunggal yang menerima \textit{event} dari seluruh \textit{event source} (R-TA-01, R-TA-03). \\
\hline
KF-TA-02 & Menangkap dan merekam keseluruhan jenis \textit{event} EV1 hingga EV13. & Format \texttt{SignedEvent} yang seragam memungkinkan ALS menerima seluruh jenis \textit{event} melalui \textit{tagged enum} \texttt{AuditEventDetails} (R-TA-03). \\
\hline
KF-TA-03 & Memformat dan menstandardisasi \textit{audit record} ke dalam skema terstruktur. & Skema \texttt{AuditRecord} bertingkat dengan atribut umum wajib dan atribut tambahan spesifik per jenis \textit{event} (R-TA-02). \\
\hline
KF-TA-04 & Menyimpan \textit{audit record} pada media penyimpanan yang bersifat \textit{append-only}. & Penulisan ke berkas log lokal menggunakan mode \textit{append} dikombinasikan dengan \textit{hash-chaining} (R-TA-04); arsip jangka panjang pada IPFS (R-TA-05). \\
\hline
KF-TA-05 & Menyediakan mekanisme retensi jangka panjang untuk mengantisipasi \textit{pruning} pada jaringan DLT. & \textit{Pinning} berkas arsip pada IPFS Cluster dan pencatatan CID beserta \textit{file\_hash} pada \textit{smart contract} Move di IOTA (R-TA-05). \\
\hline
KNF-TA-01 & Integritas \textit{audit record} yang \textit{tamper-evident}. & \textit{Hash-chaining} antar \textit{record} (R-TA-04) dan penguncian metadata berkas arsip di IOTA (R-TA-05). \\
\hline
KNF-TA-02 & Format data log yang konsisten dan mudah diproses. & Skema \texttt{AuditRecord} terstruktur dalam format JSON-Lines dengan anotasi \texttt{\#[serde(flatten)]} (R-TA-02). \\
\hline
KNF-TA-03 & Ketersediaan data audit log dalam jangka panjang. & \textit{IPFS pinning} dan penyimpanan metadata sebagai \textit{frozen object} di IOTA (R-TA-05). \\
\hline
KNF-TA-04 & Konsistensi pencatatan tanpa \textit{event} bernilai audit tinggi yang terlewat. & Verifikasi \texttt{IotaSignature} atas \texttt{EncryptedSignedEvent} pada ALS, serta \texttt{AlsQueue} dengan \textit{retry worker} pada setiap \textit{event source} untuk mengurangi risiko \textit{event} tidak terkirim (R-TA-03). \\
\hline \\
\hline

\end{xltabular}

Tabel~\ref{tab:pemenuhan-kebutuhan} menunjukkan bahwa seluruh kebutuhan fungsional (KF-TA-01 hingga KF-TA-05) dan kebutuhan nonfungsional (KNF-TA-01 hingga KNF-TA-04) terpenuhi oleh mekanisme rancangan yang telah diuraikan. Ketertelusuran dari kebutuhan ke rancangan ini menjadi dasar verifikasi pada tahap pengujian yang akan dipaparkan pada Bab IV.